\documentclass[12pt]{article}
\usepackage[utf8]{inputenc}
\usepackage[margin=1in]{geometry}
\usepackage{graphicx}
\usepackage{natbib}
\usepackage{hyperref}
\usepackage{url}

\title{Research Proposal:\\
[Data Center Energy Demand Forecasting Using LSTM]}
\author{Diego Lopez\\
Computer Science\\
University of North Carolina at Charlotte\\
Charlotte, NC, USA\\
\texttt{dlopez18@charlotte.edu}\\
\and
Matthew De La Rosa\\
Cyber Security\\
University of North Carolina at Charlotte\\
Charlotte, NC, USA\\
\texttt{mdelaro2@charlotte.edu}}
\date{\today}

\begin{document}

\maketitle

\begin{abstract}
This research proposal aims to forecast energy demand for powering data centers using machine learning techniques.
Projecting energy demand is critical for optimizing data center operations, ensuring energy efficiency, and mitigating energy-related risks.
Our goal is to develop a robust and accurate model that can accurately predict energy demand for an energy grid, enabling data centers to make informed decisions about energy usage and resource management.
\end{abstract}

\section{Introduction}
\subsection{Background and Motivation}
Recent advancements in the training of large language models like deep seek have added volitility to the projected energy consumption of generative models.
This research proposal explores the potential to use generative models to forecast energy demand for data centers. The importance of accurately forecasting energy demands
is critical for investments in energy systems to support large data center operations. This is most notable in cities that have seen explosive growth in data center construction.
Inaccurate forecasting of energy demand to support generative artificial intelligence is costly. Predicting the future consumption of energy for data centers is extremely difficult.
Instead we suggest creating a method to improve energy supply to cities with data center clusters to create more robust and scalable energy systems.

\subsection{Problem Statement}
Robust and scalable energy systems are critical for the success of data center operations.

\section{Research Objectives}
\begin{itemize}
    \item \textbf{Primary Objective:} Develop a predictive model for energy demand forecasting in data centers using reinforcement learning and binary streaming data.
    \item \textbf{Secondary Objectives:}
    \begin{itemize}
        \item Assess the impact of different machine learning algorithms on prediction accuracy and computational efficiency.
        \item Explore the integration of renewable energy sources in data center operations.
        \item Investigate the scalability of the proposed model for larger data center networks.
    \end{itemize}
    \item \textbf{Research Questions:}
    \begin{itemize}
        \item How can reinforcement learning be effectively applied to predict energy demand in data centers?
        \item What are the trade-offs between accuracy and computational cost in different modeling approaches?
        \item How can the model be adapted to incorporate renewable energy sources?
    \end{itemize}
\end{itemize}

\section{Literature Review}
Summary of relevant existing research and identification of research gaps still remain.

\section{Proposed Methodology}
\subsection{Research Design}
Our research will employ a hybrid deep learning and reinforcement learning approach to energy demand forecasting, specifically tailored for data center energy systems. The methodology will integrate the following key components:

\begin{itemize}
    \item \textbf{Deep Learning with LSTM Networks:} Utilize Long Short-Term Memory (LSTM) networks for feature extraction and modeling temporal patterns in energy consumption data.
    \item \textbf{Advanced Reinforcement Learning:} Implement a reinforcement learning framework to optimize energy allocation strategies, adapting to dynamic grid conditions.
    \item \textbf{Hybrid Model Integration:} Combine deep learning for state representation with reinforcement learning for decision-making.
\end{itemize}

\subsection{Limitations and Considerations}
\begin{itemize}
    \item \textbf{Dynamic Adaptability:} The hybrid approach allows for adaptation to changing conditions and learning from new data.
    \item \textbf{Scalability:} Deep learning models handle large datasets and complex feature spaces, improving scalability.
    \item \textbf{Real-time Optimization:} Reinforcement learning enables continuous optimization, suitable for real-time applications.
    \item \textbf{Avoid Overfitting:} Focus on capturing future trends and anomalies, not just historical data.
    \item \textbf{Incorporate External Factors:} Ensure the model accounts for external influences like weather and market dynamics.
\end{itemize}

\section{Expected Outcomes and Contributions}
More robust and scalable energy supply for cities with data center clusters.

\section{Timeline and Milestones}
To be completed during Spring 2025.

\section{Required Resources}
Access to University Research Computing (URCC) for running experiments.
Publicly Available Energy Demand Datasets from the Department of Energy and energy providers.
Publicly Available Weather and Market Dynamics Datasets.
Faculty advisor to provide guidance and support.

\section{Collaboration Plan}
The project will be led by Diego Lopez and Matthew De La Rosa.
They will work together to design and implement the proposed methodolody.
They will also collaborate on the writing of the report.

\section{Further Research Directions}
Develop secure system architecture to implement the proposed methodolody.
Develop an interface for model interpretability and explainability for stakeholders.
Validate the proposed methodolody on real-world data.


\bibliographystyle{plain}
\bibliography{references}

\end{document}
